% --- IMPORTS ---
\documentclass[leqno]{article}
\usepackage{graphicx}
\usepackage{dsfont}
\usepackage{mathtools}
\usepackage{amssymb}
\usepackage{amsmath}
\usepackage{amsthm}
\usepackage{float}
\usepackage{ragged2e}
\usepackage{tensor}
\usepackage{fancyhdr}
\usepackage{empheq}
\usepackage[most]{tcolorbox}
\usepackage{enumitem}
\usepackage{dirtytalk}
\usepackage{marginnote}
\usepackage{microtype}
\usepackage[
  a4paper,
  left=0.8in,
  right=1in,
  top=1in,
  bottom=0.5in,
  headheight=14pt,
  headsep=0.4in
]{geometry}
\setlength{\parindent}{0cm}
% --- TikZ Packages ---
\usepackage{pgfplots}
\pgfplotsset{compat=1.18}
\usepgfplotslibrary{fillbetween, polar}
\usetikzlibrary{calc, positioning}
\usetikzlibrary{angles, quotes}
\usetikzlibrary{arrows.meta}
\usetikzlibrary{decorations.pathreplacing, decorations.markings}
\usetikzlibrary{patterns}
\usetikzlibrary{intersections}
% --- URL Packages ---
\usepackage[colorlinks=true,linkcolor=red,citecolor=magenta,urlcolor=black]{hyperref}%
\urlstyle{same}
% --- END OF IMPORTS ---

% --- CUSTOM COMMANDS ---
\RenewDocumentCommand{\marks}{o m}{%
  \IfValueTF{#1}%
    {\marginnote{\raggedright\textbf{[#2]}}[#1]}%
    {\marginnote{\raggedright\textbf{[#2]}}}%
}
\newcommand{\vspacerule}[2]{%
  \par\vspace*{#1}%
  \noindent\hrulefill\par
  \vspace*{#2}%
}
\definecolor{scarlet}{RGB}{205, 40, 75}
\newcommand{\questionitem}{%
  \item\phantomsection
  \addcontentsline{toc}{section}{Question \number\value{enumi}}%
}
\renewcommand{\contentsname}{Questions}
% --- END OF CUSTOM COMMANDS ---

% --- HEADER CONFIG ---
\pagestyle{fancy}
\fancyhead{}
\fancyfoot{}
\renewcommand{\headrulewidth}{0.9pt}
\lhead{\textit{Solutions} - Complex Series}
\chead{}
\rhead{\href{https://danieljsmith.org}{danieljsmith.org}}
% --- END OF HEADER CONFIG ---

\begin{document}
\vspace*{20pt}
\tableofcontents
\newpage
\begin{enumerate}[
  label=\textbf{\arabic*.},
  leftmargin=*,
  topsep=0pt,
  partopsep=0pt,
  parsep=0pt
]

% ---- Question 1 ----
\questionitem

The infinite series $C$ and $S$ are defined by
\begin{align*}
 C&= \cos 2\theta - \frac{1}{3}\cos 5\theta + \frac{1}{9}\cos 8\theta - \frac{1}{27}\cos 11\theta + \cdots \\[4pt]
S &= \sin 2\theta - \frac{1}{3}\sin 5\theta + \frac{1}{9}\sin 8\theta - \frac{1}{27}\sin 11\theta + \cdots  
\end{align*}

Given that the series $C$ and $S$ are both convergent. \\

\begin{enumerate}[label=\textbf{(\alph*)}]
  \item Show that
  \[
  C+ \mathrm{i}S = \frac{3e^{2\mathrm{i}\theta}}{3 + e^{3\mathrm{i}\theta}}
  \]
  \marks[-20pt]{4}
  \item Hence show that
  \[
  S = \frac{9\sin 2\theta - 3\sin \theta}{10 + 6\cos 3\theta}
  \]
  \marks[-30pt]{4}
\end{enumerate}

\vspace*{10pt}

\begin{tcolorbox}[colback=white, colframe=scarlet, title={\textbf{Solution}}, breakable, grow to left by=6mm, grow to right by=10mm]
\begin{enumerate}[label=\textbf{(\alph*)}]
  \item Using Euler's formula,
\[
\cos x + \mathrm{i}\sin x = e^{\mathrm{i}x}
\]

So
\begin{align*}
C+\mathrm{i}S
&= \left(\cos 2\theta+\mathrm{i}\sin 2\theta\right)
-\frac13\left(\cos 5\theta+\mathrm{i}\sin 5\theta\right)
+\frac19\left(\cos 8\theta+\mathrm{i}\sin 8\theta\right)-\cdots \\
&= e^{2\mathrm{i}\theta}-\frac13 e^{5\mathrm{i}\theta}+\frac19 e^{8\mathrm{i}\theta}-\frac1{27}e^{11\mathrm{i}\theta}+\cdots
\end{align*}

Factor out \(e^{2\mathrm{i}\theta}\):
\begin{align*}
C+\mathrm{i}S
&= e^{2\mathrm{i}\theta}\left(1-\frac13 e^{3\mathrm{i}\theta}+\frac19 e^{6\mathrm{i}\theta}-\frac1{27}e^{9\mathrm{i}\theta}+\cdots\right)
\end{align*}

This is a geometric series with first term \(1\) and common ratio
\[
r=-\frac{e^{3\mathrm{i}\theta}}{3}
\]
Since
\[
|r|=\frac{|e^{3\mathrm{i}\theta}|}{3}=\frac13<1
\]
the series converges, so
\begin{align*}
C+\mathrm{i}S
&= e^{2\mathrm{i}\theta}\cdot \frac{1}{1-(-e^{3\mathrm{i}\theta}/3)} \\
&= e^{2\mathrm{i}\theta}\cdot \frac{1}{1+e^{3\mathrm{i}\theta}/3} \\
&= \frac{3e^{2\mathrm{i}\theta}}{3+e^{3\mathrm{i}\theta}}
\end{align*}

Hence
\[
C+\mathrm{i}S=\frac{3e^{2\mathrm{i}\theta}}{3+e^{3\mathrm{i}\theta}}
\]

  \item From part (a),
\[
C+\mathrm{i}S=\frac{3e^{2\mathrm{i}\theta}}{3+e^{3\mathrm{i}\theta}}
\]

Multiply numerator and denominator by \(3+e^{-3\mathrm{i}\theta}\):
\begin{align*}
C+\mathrm{i}S
&= \frac{3e^{2\mathrm{i}\theta}(3+e^{-3\mathrm{i}\theta})}{(3+e^{3\mathrm{i}\theta})(3+e^{-3\mathrm{i}\theta})} \\[10pt]
&= \frac{9e^{2\mathrm{i}\theta}+3e^{-\mathrm{i}\theta}}{9+3e^{3\mathrm{i}\theta}+3e^{-3\mathrm{i}\theta}+e^{3\mathrm{i}\theta}e^{-3\mathrm{i}\theta}} \\[10pt]
&= \frac{9e^{2\mathrm{i}\theta}+3e^{-\mathrm{i}\theta}}{10+3\left(e^{3\mathrm{i}\theta}+e^{-3\mathrm{i}\theta}\right)}
\end{align*}

Now use
\[
e^{3\mathrm{i}\theta}+e^{-3\mathrm{i}\theta}=2\cos 3\theta
\]
so the denominator becomes
\[
10+6\cos 3\theta
\]

Therefore
\[
C+\mathrm{i}S=\frac{9e^{2\mathrm{i}\theta}+3e^{-\mathrm{i}\theta}}{10+6\cos 3\theta}
\]

The denominator is real, so taking imaginary parts gives
\begin{align*}
S
&= \operatorname{Im}\left[\frac{9e^{2\mathrm{i}\theta}+3e^{-\mathrm{i}\theta}}{10+6\cos 3\theta}\right] \\[10pt]
&= \frac{9\sin 2\theta+3(-\sin\theta)}{10+6\cos 3\theta}
\end{align*}

Hence
\[
S=\frac{9\sin 2\theta-3\sin\theta}{10+6\cos 3\theta}
\]
\end{enumerate}
\end{tcolorbox}


\newpage

% ---- Question 2 ----
\questionitem

\begin{enumerate}[label=(\roman*)]
  \item Show that \[(1 - 2e^{2\mathrm{i}\theta})(1 - 2e^{-2\mathrm{i}\theta}) = 5 - 4\cos 2\theta\] \marks[-15pt]{3} \\
\end{enumerate}

For a positive integer $n$, series $C$ and $S$ are defined by\\

\[
\begin{aligned}
C &= \cos\theta + 2\cos 3\theta + 4\cos 5\theta + \cdots + 2^{n-1}\cos((2n-1)\theta)\\[5pt]
S &= \sin\theta + 2\sin 3\theta + 4\sin 5\theta + \cdots + 2^{n-1}\sin((2n-1)\theta)
\end{aligned}
\] \\
\begin{enumerate}[label=(\roman*), resume]
  \item Show that
  \[
  C = \frac{2^{n+1}\cos((2n-1)\theta) - 2^n\cos((2n+1)\theta) - \cos\theta}{5 - 4\cos 2\theta}
  \]
  and find a similar expression for $S$. \marks{9}
\end{enumerate}

\vspace*{10pt}

\begin{tcolorbox}[colback=white, colframe=scarlet, title={\textbf{Solution}}, breakable, grow to left by=6mm, grow to right by=10mm]
\begin{enumerate}[label=\textbf{(\roman*)}]
  \item Expanding the product,
  \begin{align*}
  (1-2e^{2\mathrm{i}\theta})(1-2e^{-2\mathrm{i}\theta})
  &=1-2e^{2\mathrm{i}\theta}-2e^{-2\mathrm{i}\theta}+4\\
  &=5-2\left(e^{2\mathrm{i}\theta}+e^{-2\mathrm{i}\theta}\right)
  \end{align*}
  Using
  \[
  e^{2\mathrm{i}\theta}+e^{-2\mathrm{i}\theta}=2\cos 2\theta
  \]
  gives
  \begin{align*}
  (1-2e^{2\mathrm{i}\theta})(1-2e^{-2\mathrm{i}\theta})
  &=5-2(2\cos 2\theta)\\
  &=5-4\cos 2\theta
  \end{align*}
  as required.

  \item Let
  \[
  Z=C+\mathrm{i}S
  \]
  Then, using $\cos x+\mathrm{i}\sin x=e^{\mathrm{i}x}$,
  \begin{align*}
  Z&=(\cos\theta+\mathrm{i}\sin\theta)+2(\cos3\theta+\mathrm{i}\sin3\theta)+\cdots+2^{n-1}(\cos((2n-1)\theta)+\mathrm{i}\sin((2n-1)\theta))\\
  &=e^{\mathrm{i}\theta}+2e^{3\mathrm{i}\theta}+4e^{5\mathrm{i}\theta}+\cdots+2^{n-1}e^{(2n-1)\mathrm{i}\theta}\\
  &=e^{\mathrm{i}\theta}\left(1+2e^{2\mathrm{i}\theta}+(2e^{2\mathrm{i}\theta})^2+\cdots+(2e^{2\mathrm{i}\theta})^{n-1}\right)
  \end{align*}
  This is a geometric series with first term $1$ and common ratio $2e^{2\mathrm{i}\theta}$, so
  \[
  Z=e^{\mathrm{i}\theta}\frac{1-(2e^{2\mathrm{i}\theta})^n}{1-2e^{2\mathrm{i}\theta}}
  =e^{\mathrm{i}\theta}\frac{1-2^n e^{2n\mathrm{i}\theta}}{1-2e^{2\mathrm{i}\theta}}
  \]
  Multiply numerator and denominator by $1-2e^{-2\mathrm{i}\theta}$:
  \begin{align*}
  Z&=e^{\mathrm{i}\theta}\frac{1-2^n e^{2n\mathrm{i}\theta}}{1-2e^{2\mathrm{i}\theta}}\cdot\frac{1-2e^{-2\mathrm{i}\theta}}{1-2e^{-2\mathrm{i}\theta}}\\
  &=\frac{e^{\mathrm{i}\theta}(1-2^n e^{2n\mathrm{i}\theta})(1-2e^{-2\mathrm{i}\theta})}{(1-2e^{2\mathrm{i}\theta})(1-2e^{-2\mathrm{i}\theta})}
  \end{align*}
  Using part (i), the denominator is $5-4\cos2\theta$, so
  \begin{align*}
  Z&=\frac{e^{\mathrm{i}\theta}(1-2^n e^{2n\mathrm{i}\theta})(1-2e^{-2\mathrm{i}\theta})}{5-4\cos2\theta}\\[10pt]
  &=\frac{e^{\mathrm{i}\theta}-2e^{-\mathrm{i}\theta}-2^n e^{\mathrm{i}(2n+1)\theta}+2^{n+1}e^{\mathrm{i}(2n-1)\theta}}{5-4\cos2\theta}
  \end{align*}
  Now simplify the first two terms:
  \begin{align*}
  e^{\mathrm{i}\theta}-2e^{-\mathrm{i}\theta}
  &=(\cos\theta+\mathrm{i}\sin\theta)-2(\cos\theta-\mathrm{i}\sin\theta)\\
  &=-\cos\theta+3\mathrm{i}\sin\theta
  \end{align*}
  Hence
  \[
  Z=\frac{-\cos\theta+3\mathrm{i}\sin\theta-2^n e^{\mathrm{i}(2n+1)\theta}+2^{n+1}e^{\mathrm{i}(2n-1)\theta}}{5-4\cos2\theta}
  \]
  Writing the exponential terms in sine and cosine form,
  \begin{align*}
  Z
  &=\frac{-\cos\theta-2^n\cos((2n+1)\theta)+2^{n+1}\cos((2n-1)\theta)}{5-4\cos2\theta}\\[10pt]
  &\qquad+\mathrm{i}\,\frac{3\sin\theta-2^n\sin((2n+1)\theta)+2^{n+1}\sin((2n-1)\theta)}{5-4\cos2\theta}
  \end{align*}
  Since $Z=C+\mathrm{i}S$, equating real and imaginary parts gives
  \[
  C=\frac{2^{n+1}\cos((2n-1)\theta)-2^n\cos((2n+1)\theta)-\cos\theta}{5-4\cos2\theta}\\[10pt]
  \]
  and
  \[
  S=\frac{2^{n+1}\sin((2n-1)\theta)-2^n\sin((2n+1)\theta)+3\sin\theta}{5-4\cos2\theta}\\[10pt]
  \]
  So the required similar expression for $S$ is
  \[
  S=\frac{2^{n+1}\sin((2n-1)\theta)-2^n\sin((2n+1)\theta)+3\sin\theta}{5-4\cos2\theta}\\[10pt]
  \]
\end{enumerate}
\end{tcolorbox}


\newpage

% ---- Question 3 ----
\questionitem

\begin{enumerate}[label=\textbf{(\roman*)}]
  \item Show that
  \[
  1 - e^{i2\theta} = 2\sin\theta(\sin\theta - i\cos\theta)
  \] \marks[-15pt]{3}\\
  \item For a positive integer $n$, the series $C$ and $S$ are defined as follows.
  \begin{align*}
  C &= 1 - \binom{n}{1}\cos 2\theta + \binom{n}{2}\cos 4\theta - \ldots + (-1)^n\cos 2n\theta \\[5pt]
  S &=\quad -\binom{n}{1}\sin 2\theta + \binom{n}{2}\sin 4\theta - \ldots + (-1)^n\sin 2n\theta
  \end{align*} \\
  By considering $C + iS$, show that
  \[
  C = 2^n\sin^n\theta\cos\left(n\theta - \frac{n\pi}{2}\right)
  \]
  and find a corresponding expression for $S$.\marks{9} \\
\end{enumerate}

\vspace*{10pt}

\begin{tcolorbox}[colback=white, colframe=scarlet, title={\textbf{Solution}}, breakable, grow to left by=6mm, grow to right by=10mm]
\begin{enumerate}[label=\textbf{(\roman*)}]
  \item Using Euler's formula \(e^{i2\theta}=\cos 2\theta+i\sin 2\theta\),
\begin{align*}
1-e^{i2\theta}
&=1-\cos 2\theta-i\sin 2\theta\\
&=(1-\cos 2\theta)-i\sin 2\theta\\
&=2\sin^2\theta-i(2\sin\theta\cos\theta)\\
&=2\sin\theta(\sin\theta-i\cos\theta)
\end{align*}
Hence
\[
1-e^{i2\theta}=2\sin\theta(\sin\theta-i\cos\theta)
\]

  \item Combine the two series into one complex expression:
\begin{align*}
C+iS
&=1-\binom{n}{1}(\cos 2\theta+i\sin 2\theta)+\binom{n}{2}(\cos 4\theta+i\sin 4\theta)-\cdots+(-1)^n(\cos 2n\theta+i\sin 2n\theta)\\
&=1-\binom{n}{1}e^{i2\theta}+\binom{n}{2}e^{i4\theta}-\cdots+(-1)^n e^{i2n\theta}\\
&=\sum_{r=0}^{n}\binom{n}{r}(-1)^r e^{i2r\theta}\\
&=\sum_{r=0}^{n}\binom{n}{r}\bigl(-e^{i2\theta}\bigr)^r\\
&=(1-e^{i2\theta})^n
\end{align*}
by the binomial theorem.

From part (i),
\[
1-e^{i2\theta}=2\sin\theta(\sin\theta-i\cos\theta)
\]
so
\begin{align*}
C+iS
&=\left[2\sin\theta(\sin\theta-i\cos\theta)\right]^n\\
&=2^n\sin^n\theta\,(\sin\theta-i\cos\theta)^n
\end{align*}

Now factor out \(-i\):
\[
\sin\theta-i\cos\theta
=
-i(\cos\theta+i\sin\theta)
\]
Since
\[
-i=e^{-i\pi/2}
\]
this becomes
\[
\sin\theta-i\cos\theta
=
e^{-i\pi/2}(\cos\theta+i\sin\theta)
\]

Therefore, by De Moivre's theorem,
\begin{align*}
(\sin\theta-i\cos\theta)^n
&=\left(e^{-i\pi/2}(\cos\theta+i\sin\theta)\right)^n\\
&=e^{-in\pi/2}\,e^{in\theta}\\
&=\cos\left(n\theta-\frac{n\pi}{2}\right)+i\sin\left(n\theta-\frac{n\pi}{2}\right)
\end{align*}

Substituting this into \(C+iS\),
\begin{align*}
C+iS
&=2^n\sin^n\theta\left(\cos\left(n\theta-\frac{n\pi}{2}\right)+i\sin\left(n\theta-\frac{n\pi}{2}\right)\right)
\end{align*}

Equating real and imaginary parts gives
\[
C=2^n\sin^n\theta\cos\left(n\theta-\frac{n\pi}{2}\right)
\]
and
\[
S=2^n\sin^n\theta\sin\left(n\theta-\frac{n\pi}{2}\right)
\]
\end{enumerate}
\end{tcolorbox}


\newpage

% ---- Question 4 ----
\questionitem

The infinite series $C$ and $S$ are defined as follows: \\
\begin{align*}
C&= \cos\alpha + r\cos(\alpha + \beta) + r^2\cos(\alpha + 2\beta) + \ldots \\[10pt]
S &= \sin\alpha + r\sin(\alpha + \beta) + r^2\sin(\alpha + 2\beta) + \ldots \\
\end{align*}

where $r$ is a real number and $|r| < 1$. \\


By considering $C+ \mathrm{i}S$, show that \\
\[
C= \frac{\cos\alpha - r\cos(\alpha - \beta)}{1 + r^2 - 2r\cos\beta} \\
\] 
and find a corresponding expression for $S$. \marks{8}

\vspace*{10pt}

\begin{tcolorbox}[colback=white, colframe=scarlet, title={\textbf{Solution}}, breakable, grow to left by=6mm, grow to right by=10mm]
Let
\[
Z=C+\mathrm{i}S
\]

Then
\begin{align*}
Z&=\left(\cos\alpha+r\cos(\alpha+\beta)+r^2\cos(\alpha+2\beta)+\cdots\right) \\
&\qquad +\mathrm{i}\left(\sin\alpha+r\sin(\alpha+\beta)+r^2\sin(\alpha+2\beta)+\cdots\right) \\
&=\sum_{n=0}^{\infty} r^n\left(\cos(\alpha+n\beta)+\mathrm{i}\sin(\alpha+n\beta)\right)
\end{align*}

Using Euler's formula,
\[
\cos\theta+\mathrm{i}\sin\theta=e^{\mathrm{i}\theta}
\]

so
\begin{align*}
Z&=\sum_{n=0}^{\infty} r^n e^{\mathrm{i}(\alpha+n\beta)} \\
&=e^{\mathrm{i}\alpha}\sum_{n=0}^{\infty}\left(re^{\mathrm{i}\beta}\right)^n
\end{align*}

Since
\[
\left|re^{\mathrm{i}\beta}\right|=|r|<1
\]
this is a convergent geometric series, so
\begin{align*}
Z&=e^{\mathrm{i}\alpha}\cdot \frac{1}{1-re^{\mathrm{i}\beta}} \\
&=\frac{e^{\mathrm{i}\alpha}}{1-re^{\mathrm{i}\beta}}
\end{align*}

Now multiply top and bottom by the conjugate-like factor \(1-re^{-\mathrm{i}\beta}\):
\begin{align*}
Z&=\frac{e^{\mathrm{i}\alpha}(1-re^{-\mathrm{i}\beta})}{(1-re^{\mathrm{i}\beta})(1-re^{-\mathrm{i}\beta})}
\end{align*}

First simplify the denominator:
\begin{align*}
(1-re^{\mathrm{i}\beta})(1-re^{-\mathrm{i}\beta})
&=1-r(e^{\mathrm{i}\beta}+e^{-\mathrm{i}\beta})+r^2 \\
&=1-r(2\cos\beta)+r^2 \\
&=1+r^2-2r\cos\beta
\end{align*}

Now simplify the numerator:
\begin{align*}
e^{\mathrm{i}\alpha}(1-re^{-\mathrm{i}\beta})
&=e^{\mathrm{i}\alpha}-re^{\mathrm{i}\alpha}e^{-\mathrm{i}\beta} \\
&=e^{\mathrm{i}\alpha}-re^{\mathrm{i}(\alpha-\beta)}
\end{align*}

Hence
\[
Z=\frac{e^{\mathrm{i}\alpha}-re^{\mathrm{i}(\alpha-\beta)}}{1+r^2-2r\cos\beta}
\]

Write the exponentials in terms of sine and cosine:
\begin{align*}
Z&=\frac{\cos\alpha+\mathrm{i}\sin\alpha-r\left(\cos(\alpha-\beta)+\mathrm{i}\sin(\alpha-\beta)\right)}{1+r^2-2r\cos\beta} \\
&=\frac{\left(\cos\alpha-r\cos(\alpha-\beta)\right)+\mathrm{i}\left(\sin\alpha-r\sin(\alpha-\beta)\right)}{1+r^2-2r\cos\beta}
\end{align*}

But \(Z=C+\mathrm{i}S\), so equating real and imaginary parts gives
\[
C=\frac{\cos\alpha-r\cos(\alpha-\beta)}{1+r^2-2r\cos\beta}
\]
and
\[
S=\frac{\sin\alpha-r\sin(\alpha-\beta)}{1+r^2-2r\cos\beta}
\]

Therefore,
\[
C=\frac{\cos\alpha-r\cos(\alpha-\beta)}{1+r^2-2r\cos\beta},
\qquad
S=\frac{\sin\alpha-r\sin(\alpha-\beta)}{1+r^2-2r\cos\beta}
\]
\end{tcolorbox}


\newpage

% ---- Question 5 ----
\questionitem

\begin{enumerate}[label=(\roman*)]
  \item For values of $\theta$ for which the expressions are defined, show that
  \[
  1+\mathrm{i}\tan\theta=\sec\theta(\cos\theta+\mathrm{i}\sin\theta)
  \]
  \marks[-20pt]{2}

  \item For a positive integer $n$, the series $C$ and $S$ are defined by
  \begin{align*}
    C&=\,1-\binom{n}{2}\tan^2\theta+\binom{n}{4}\tan^4\theta-\cdots\\[10pt]
  S&=\qquad\binom{n}{1}\tan\theta-\binom{n}{3}\tan^3\theta+\binom{n}{5}\tan^5\theta-\cdots
  \end{align*}

  where in each case the pattern continues until the power of $\tan\theta$ is at most $n$. \\

  By considering $C+\mathrm{i}S$, show that
  \[
  C=\sec^n\theta\cos n\theta
  \]
  and find a corresponding expression for $S$. \marks{7}
\end{enumerate}

\vspace*{10pt}

\begin{tcolorbox}[colback=white, colframe=scarlet, title={\textbf{Solution}}, breakable, grow to left by=6mm, grow to right by=10mm]
\begin{enumerate}[label=\textbf{(\roman*)}]
\item For the expressions to be defined, we need
\[
\cos\theta\neq 0
\]

Now
\begin{align*}
\sec\theta(\cos\theta+\mathrm{i}\sin\theta)
&=\sec\theta\cos\theta+\mathrm{i}\sec\theta\sin\theta \\
&=1+\mathrm{i}\tan\theta
\end{align*}

So
\[
1+\mathrm{i}\tan\theta=\sec\theta(\cos\theta+\mathrm{i}\sin\theta)
\]

as required.

\item Consider \(C+\mathrm{i}S\).

Using the given definitions,
\begin{align*}
C+\mathrm{i}S
&=\left(1-\binom{n}{2}\tan^2\theta+\binom{n}{4}\tan^4\theta-\cdots\right)
+\mathrm{i}\left(\binom{n}{1}\tan\theta-\binom{n}{3}\tan^3\theta+\binom{n}{5}\tan^5\theta-\cdots\right)
\end{align*}

Now note that
\[
(\mathrm{i}\tan\theta)^0=1,\qquad
(\mathrm{i}\tan\theta)^1=\mathrm{i}\tan\theta,\qquad
(\mathrm{i}\tan\theta)^2=-\tan^2\theta,\qquad
(\mathrm{i}\tan\theta)^3=-\mathrm{i}\tan^3\theta
\]

so the terms in \(C+\mathrm{i}S\) match the binomial expansion of \((1+\mathrm{i}\tan\theta)^n\). Hence
\[
C+\mathrm{i}S=\sum_{k=0}^{n}\binom{n}{k}(\mathrm{i}\tan\theta)^k=(1+\mathrm{i}\tan\theta)^n
\]

Using part (i),
\[
1+\mathrm{i}\tan\theta=\sec\theta(\cos\theta+\mathrm{i}\sin\theta)
\]

Therefore
\begin{align*}
C+\mathrm{i}S
&=(1+\mathrm{i}\tan\theta)^n \\
&=\left[\sec\theta(\cos\theta+\mathrm{i}\sin\theta)\right]^n \\
&=\sec^n\theta\,(\cos\theta+\mathrm{i}\sin\theta)^n
\end{align*}

By de Moivre's theorem,
\[
(\cos\theta+\mathrm{i}\sin\theta)^n=\cos n\theta+\mathrm{i}\sin n\theta
\]

So
\[
C+\mathrm{i}S=\sec^n\theta(\cos n\theta+\mathrm{i}\sin n\theta)
\]

Equating real and imaginary parts gives
\[
C=\sec^n\theta\cos n\theta
\]
and
\[
S=\sec^n\theta\sin n\theta
\]

Hence the corresponding expression for \(S\) is
\[
S=\sec^n\theta\sin n\theta
\]
\end{enumerate}
\end{tcolorbox}


\newpage

% ---- Question 6 ----
\questionitem

The convergent series $C$ and $S$ are defined by
\begin{align*}
C &= \sum_{n=0}^{\infty} \left(-\frac{1}{2}\right)^n \cos\bigl((4n+1)\theta\bigr)
\\[10pt]
S &= \sum_{n=0}^{\infty} \left(-\frac{1}{2}\right)^n \sin\bigl((4n+1)\theta\bigr)
\end{align*} \\
\begin{enumerate}[label=\textbf{(\alph*)}]
  \item Show that
  \[
  C + \mathrm{i}S = \frac{2e^{\mathrm{i}\theta}}{2 + e^{4\mathrm{i}\theta}}
  \]
  \marks[-20pt]{4}
  \item Hence show that
  \[
  C = \frac{4\cos \theta + 2\cos 3\theta}{5 + 4\cos 4\theta}
  \]
  \marks[-30pt]{4}
\end{enumerate}

\vspace*{10pt}

\begin{tcolorbox}[colback=white, colframe=scarlet, title={\textbf{Solution}}, breakable, grow to left by=6mm, grow to right by=10mm]
\begin{enumerate}[label=\textbf{(\alph*)}]
  \item Using Euler's formula,
\[
\cos x+\mathrm{i}\sin x=e^{\mathrm{i}x}
\]
so
\begin{align*}
C+\mathrm{i}S
&=\sum_{n=0}^{\infty}\left(-\frac12\right)^n\left(\cos\bigl((4n+1)\theta\bigr)+\mathrm{i}\sin\bigl((4n+1)\theta\bigr)\right)\\
&=\sum_{n=0}^{\infty}\left(-\frac12\right)^n e^{\mathrm{i}(4n+1)\theta}\\
&=e^{\mathrm{i}\theta}\sum_{n=0}^{\infty}\left(-\frac{e^{4\mathrm{i}\theta}}{2}\right)^n
\end{align*}

This is a geometric series with common ratio
\[
r=-\frac{e^{4\mathrm{i}\theta}}{2}
\]
and
\[
|r|=\frac12<1
\]
so it converges and its sum is
\begin{align*}
C+\mathrm{i}S
&=e^{\mathrm{i}\theta}\cdot \frac{1}{1-r}\\
&=e^{\mathrm{i}\theta}\cdot \frac{1}{1+\frac12 e^{4\mathrm{i}\theta}}\\
&=\frac{2e^{\mathrm{i}\theta}}{2+e^{4\mathrm{i}\theta}}
\end{align*}

Hence
\[
C+\mathrm{i}S=\frac{2e^{\mathrm{i}\theta}}{2+e^{4\mathrm{i}\theta}}
\]

  \item Starting from part (a),
\[
C+\mathrm{i}S=\frac{2e^{\mathrm{i}\theta}}{2+e^{4\mathrm{i}\theta}}
\]

Multiply top and bottom by \(2+e^{-4\mathrm{i}\theta}\):
\begin{align*}
C+\mathrm{i}S
&=\frac{2e^{\mathrm{i}\theta}(2+e^{-4\mathrm{i}\theta})}{(2+e^{4\mathrm{i}\theta})(2+e^{-4\mathrm{i}\theta})}\\
&=\frac{4e^{\mathrm{i}\theta}+2e^{-3\mathrm{i}\theta}}{(2+e^{4\mathrm{i}\theta})(2+e^{-4\mathrm{i}\theta})}
\end{align*}

Now simplify the denominator:
\begin{align*}
(2+e^{4\mathrm{i}\theta})(2+e^{-4\mathrm{i}\theta})
&=4+2e^{4\mathrm{i}\theta}+2e^{-4\mathrm{i}\theta}+1\\
&=5+2\left(e^{4\mathrm{i}\theta}+e^{-4\mathrm{i}\theta}\right)\\
&=5+4\cos 4\theta
\end{align*}
since \(e^{\mathrm{i}x}+e^{-\mathrm{i}x}=2\cos x\).

Therefore
\[
C+\mathrm{i}S=\frac{4e^{\mathrm{i}\theta}+2e^{-3\mathrm{i}\theta}}{5+4\cos 4\theta}
\]

The denominator is real, so taking real parts gives
\begin{align*}
C
&=\frac{4\cos\theta+2\cos(-3\theta)}{5+4\cos4\theta}\\[10pt]
&=\frac{4\cos\theta+2\cos3\theta}{5+4\cos4\theta}
\end{align*}

Hence
\[
C=\frac{4\cos \theta + 2\cos 3\theta}{5 + 4\cos 4\theta}
\]
\end{enumerate}
\end{tcolorbox}


\newpage

% ---- Question 7 ----
\questionitem

Let
\begin{align*}
C&= \sum_{r=0}^{n-1}\cos^2\left(\theta + \frac{2\pi r}{n}\right)\\[10pt]
S &= \sum_{r=0}^{n-1}\sin\left(\theta + \frac{2\pi r}{n}\right)\cos\left(\theta + \frac{2\pi r}{n}\right)\\
\end{align*}

where $n$ is an integer greater than $2$. Let $\exp(x)$ denote $e^x$.\\


By considering the real and imaginary parts of
\[
\sum_{r=0}^{n-1} \exp\left(2\mathrm{i}\left(\theta + \dfrac{2\pi r}{n}\right)\right)
\]
show that $C= \dfrac{n}{2}$ and $S = 0$. \marks{7} \\

\vspace*{10pt}

\begin{tcolorbox}[colback=white, colframe=scarlet, title={\textbf{Solution}}, breakable, grow to left by=6mm, grow to right by=10mm]
Let
\[
Z=\sum_{r=0}^{n-1}\exp\left(2\mathrm{i}\left(\theta+\frac{2\pi r}{n}\right)\right)
\]

Using \(\exp(\mathrm{i}x)=\cos x+\mathrm{i}\sin x\), this becomes
\begin{align*}
Z
&=\sum_{r=0}^{n-1}\left(\cos\left(2\theta+\frac{4\pi r}{n}\right)+\mathrm{i}\sin\left(2\theta+\frac{4\pi r}{n}\right)\right)
\end{align*}

So
\begin{align*}
\operatorname{Re}(Z)&=\sum_{r=0}^{n-1}\cos\left(2\theta+\frac{4\pi r}{n}\right) \\
\operatorname{Im}(Z)&=\sum_{r=0}^{n-1}\sin\left(2\theta+\frac{4\pi r}{n}\right)
\end{align*}

Now rewrite \(Z\) as a geometric series:
\begin{align*}
Z
&=\sum_{r=0}^{n-1}\exp\left(2\mathrm{i}\theta+\frac{4\pi \mathrm{i}r}{n}\right) \\
&=\exp(2\mathrm{i}\theta)\sum_{r=0}^{n-1}\left(\exp\left(\frac{4\pi \mathrm{i}}{n}\right)\right)^r
\end{align*}

Let
\[
\omega=\exp\left(\frac{4\pi \mathrm{i}}{n}\right)
\]

Then
\[
Z=\exp(2\mathrm{i}\theta)\sum_{r=0}^{n-1}\omega^r
\]

Since \(n>2\), we have \(0<\frac{2}{n}<1\), so \(\omega\neq 1\). Therefore the geometric series formula gives
\begin{align*}
Z
&=\exp(2\mathrm{i}\theta)\frac{1-\omega^n}{1-\omega}
\end{align*}

But
\begin{align*}
\omega^n
&=\left(\exp\left(\frac{4\pi \mathrm{i}}{n}\right)\right)^n \\[5pt]
&=\exp(4\pi \mathrm{i}) \\
&=1
\end{align*}

So
\begin{align*}
Z
&=\exp(2\mathrm{i}\theta)\frac{1-1}{1-\omega} \\
&=0
\end{align*}

Hence both the real and imaginary parts of \(Z\) are zero:
\begin{align*}
\sum_{r=0}^{n-1}\cos\left(2\theta+\frac{4\pi r}{n}\right)&=0 \\
\sum_{r=0}^{n-1}\sin\left(2\theta+\frac{4\pi r}{n}\right)&=0
\end{align*}

Now use the double-angle identities.

For \(C\),
\[
\cos^2 x=\frac{1+\cos 2x}{2}
\]

with \(x=\theta+\frac{2\pi r}{n}\). Then
\begin{align*}
C
&=\sum_{r=0}^{n-1}\cos^2\left(\theta+\frac{2\pi r}{n}\right) \\[10pt]
&=\sum_{r=0}^{n-1}\frac{1+\cos\left(2\theta+\frac{4\pi r}{n}\right)}{2} \\[10pt]
&=\frac12\sum_{r=0}^{n-1}1+\frac12\sum_{r=0}^{n-1}\cos\left(2\theta+\frac{4\pi r}{n}\right) \\[10pt]
&=\frac12(n)+\frac12(0) =\frac{n}{2}
\end{align*}

For \(S\),
\[
\sin x\cos x=\frac12\sin 2x
\]

so, again with \(x=\theta+\frac{2\pi r}{n}\),
\begin{align*}
S
&=\sum_{r=0}^{n-1}\sin\left(\theta+\frac{2\pi r}{n}\right)\cos\left(\theta+\frac{2\pi r}{n}\right) \\[10pt]
&=\frac12\sum_{r=0}^{n-1}\sin\left(2\theta+\frac{4\pi r}{n}\right) \\[10pt]
&=\frac12(0) 
=0
\end{align*}

Therefore
\[
C=\frac{n}{2}
\qquad \text{and} \qquad
S=0
\]
\end{tcolorbox}


\newpage
\end{enumerate}
\end{document}
